
\chapter{Maria i tempelet}

Da Maria fylte to år, kom Joakim til Anna en kveld.
Solen var på vei ned, og de satt sammen i gårdsplassen.
Skyggene fra oliventrærne ble lange over steinene.

Joakim var stille lenge før han begynte å snakke.
«Anna.» Hun så opp. «Jeg har tenkt mye.»
«På hva?», spurte hun.
Han så mot huset der Maria lekte.
«På løftet vi ga.»

Anna fulgte blikket hans.
Hun så datteren løpe etter en sommerfugl.
Et smil kom over ansiktet hennes.
«Ja,» sa hun lavt.

Joakim sukket.
«Vi lovet Herren at når tiden kom, skulle vi føre henne til tempelet.»
Ordene hang i luften.
Anna kjente et stikk i hjertet.
Hun hadde visst at denne dagen ville komme.
Men hun hadde likevel håpet at den kunne vente.
«Jeg vet», svarte hun.

Joakim tok hånden hennes og sa:
«Vi fikk henne som en gave. Vi må holde ord.
Jeg ønsker ikke at vårt løfte skal bli tomme ord foran Gud.»
Anna så ned og svarte:
«Tror du Herren ville forkaste oss fordi vi ønsker litt mer tid med henne?»
«Nei», svart han med et trist smil.
«Men jeg ønsker å gjøre det rette.»

Anna var stille. Så sa hun:
«La oss vente ett år.»
Joakim så på henne. «Ett år?»
Hun nikket.
«La henne bli tre år. Da vil hun være sterkere.
Da vil hun kanskje forstå mer. Og kanskje vil hun ikke
gråte etter oss når hun går.»
Joakim så lenge på henne.
Så smilte han svakt og sa:
«Du er fortsatt hennes mor.»
Anna lo stille.
«Og du er fortsatt hennes far.»
Han nikket.
«Da venter vi.»

Det året gikk raskere enn de ønsket.
Snart kom dagen da Maria fylte tre år.
Hele huset var fylt av både glede og sorg.
For alle visste at dette ikke bare var en bursdag.
Det var dagen da et gammelt løfte skulle oppfylles.

Joakim stod tidlig opp.
Han gikk ut i hagen.
Han så på treet der Anna hadde bedt sin første bønn om et barn.
Bladene på laurbærtreet beveget seg svakt i vinden.
Han la hånden på stammen.
«Her begynte alt», hvisket han.

Da reisen til tempelet skulle begynne,
kalte Anna sammen unge jenter fra familiene i Israel.
De var rene og respektfulle, og de fikk en spesiell oppgave.
De skulle bære lamper foran Maria.
Små oljelamper med flammer som skulle lyse veien.

«Hvorfor skal de ha lys?» spurte Maria.
Anna knelte foran henne og festet klærne hennes.

«Fordi veien du går, er en spesiell vei.»
Maria så alvorlig på henne.

«Skal jeg ikke komme hjem igjen?» spurte Maria.
Spørsmålet traff Anna hardt.
Hun smilte likevel.
«Du vil alltid ha et hjem i hjertet mitt.»
Hun kysset pannen hennes.
«Og Gud vil alltid være nær deg.»

De begynte reisen mot Jerusalem.
Maria gikk mellom de unge jentene som bar lampene.
Flammene skinte i dagslyset.
Hun snudde seg én gang.
Bare én.
Hun så tilbake mot foreldrene.

Anna holdt pusten.
Men Maria stoppet ikke.
Hun vendte seg fram igjen og fortsatte.

Joakim så på Anna. «Hun er klar.»
Anna tørket en tåre, og sa: «Hun er større enn jeg trodde.»

Da de kom til tempelet, stod de store steinmurene foran dem.
Sollyset skinte på de hellige bygningene.
Prester gikk gjennom gårdsplassene.
Røkelsen steg opp mot himmelen.
Maria så seg rundt med undring.
Hun var liten blant de store søylene.
Men hun var ikke redd.
En prest kom mot dem.

Han hadde et alvorlig ansikt, men øynene hans var milde.
Han bøyde seg ned til Maria.
Han kysset henne på pannen. Så sa han:
«Herren Gud har latt ditt navn bli stort blant slektene.»

Han så mot Anna og Joakim.
«Gjennom dette barnet vil Gud en dag vise sin barmhjertighet mot sitt folk.»
Alle ble stille.
Ordene fylte rommet.
Presten tok Maria med seg.
Han førte henne fram mot alteret.
Der satte han henne på det tredje trinnet.

Et lite barn.
Midt i tempelets storhet.
Men da hun stod der, skjedde noe underlig.
Det var som om en fred fylte stedet.
Som om hun hørte hjemme der.
Maria smilte.
Så begynte hun å danse.
Ikke som et barn som leker.
Men med en glede som kom fra et sted dypere enn ord.
Folkene som så det, ble rørt.
De smilte.

De hvisket til hverandre.
«Se det barnet. Hun er velsignet.»

Og fra den dagen bar Israels folk kjærlighet til
den lille jenta som hadde kommet til tempelet
med en mors tårer og et løfte som endelig var blitt oppfylt.

Anna stod et stykke unna.
Hjertet hennes var både knust og fylt av glede.
Hun visste at hun hadde gitt Maria til Gud.
Men hun visste også noe annet:
Noen gaver blir ikke mindre fordi man deler dem.
Noen gaver vokser.


Joakim og Anna gikk ned fra tempelhøyden med hjertene fulle av undring.
De hadde overlatt sin eneste datter til Herren, og selv om sorgen
stakk i dem da de forlot henne, var den blandet med en dyp glede.
De lovpriste Gud for det de hadde fått se: Maria hadde ikke nølt,
ikke løpt tilbake mot dem eller grått etter foreldrene sine.
Med barnlig tillit hadde hun gått inn i Herrens hus som om hun
allerede hørte hjemme der.

Fra den dagen levde Maria i tempelet. Prestene og kvinnene
som tjente ved helligdommen tok seg av henne, men det ble snart
sagt at det var noe uvanlig ved barnet. Hun vokste opp stille og
from, ren av sinn og mild av vesen. Hun beveget seg gjennom
templets forgårder med samme fred som en due, og alle som så henne,
la merke til den roen hun bar med seg. Tradisjonen fortalte at
Herren selv sørget for henne, og at en engel kom til
henne med mat, slik at hun levde under Guds særlige omsorg.

Årene gikk, og Maria ble tolv år gammel.
Da samlet prestene seg til råd. De visste at hun ikke kunne bli
værende i tempelet slik et lite barn kunne. Loven satte grenser,
og de fryktet at de, om de handlet feil, kunne vanhellige Herrens helligdom.

«Maria har bodd i Herrens tempel i tolv år», sa de til hverandre.
«Hva skal vi gjøre med henne? Vi kan ikke la noe skje som
gjør helligdommen uren.»

Til slutt vendte de seg til øverstepresten, Sakarja.

«Du som gjør tjeneste ved Herrens alter, gå inn og søk Guds vilje.
Be for henne. Hvis Herren åpenbarer hva som skal gjøres, vil vi følge det.»

Øverstepresten tok på seg den hellige prestedrakten med de
tolv små bjellene langs kanten, slik den skulle bæres når han
trådte inn for Herrens ansikt. Med ærefrykt gikk han gjennom
forhenget og inn i Det aller helligste,
stedet ingen kunne gå inn uten den største varsomhet.
Der ba han om visdom, slik at Gud selv skulle
vise hvilken vei Maria skulle følge.

Mens han ba, stod plutselig en Herrens engel ved siden av ham.
«Sakarja, Sakarja», sa engelen,
«gå ut og kall sammen alle enkemennene blant Israels folk.
La hver av dem komme med en stav i hånden.
Den mannen Herren utpeker ved et tegn, skal få Maria til hustru.»
Sakarja bøyde hodet og tok imot budskapet.

Straks ble herolder sendt ut i hele Judea.
De vandret fra landsby til landsby og lot
Herrens trompet lyde over åser og daler.
Kallet gikk ut til alle enkemenn i landet om å møte
frem ved tempelet med stavene sine.

Da trompettonene lød gjennom landet,
skyndte mennene seg av sted. Fra byer, landsbyer og
enslige gårder kom de vandrende mot Jerusalem,
hver med sin stav i hånden, uten å vite hvorfor Herren hadde kalt dem sammen.
